\documentclass[12pt,a4paper]{article}
\usepackage[UTF8]{ctex}
\usepackage{amsmath,amssymb,amsthm}
\usepackage{geometry}
\usepackage{bm}
\usepackage{hyperref}
\usepackage{tikz}
\usetikzlibrary{arrows,arrows.meta,positioning,calc,shapes.geometric}

\geometry{margin=2.5cm}

\title{12.1.6 平面图形方法的详细理解}
\author{第12章抓取与操作补充说明}
\date{}

\newtheorem{theorem}{定理}
\newtheorem{definition}{定义}
\newtheorem{example}{例子}

\begin{document}

\maketitle

\section{章节概述}

12.1.6章节"平面图形方法"（Planar Graphical Methods）介绍了一种在\textbf{平面问题}中可视化可行运动的图形方法。由于平面twist空间是3维的，我们可以用图形方法来表示和分析约束，这对于理解接触运动学非常有帮助。

\textbf{核心思想}：将平面twist表示为\textbf{旋转中心}（Center of Rotation, CoR），从而在2D平面上可视化3D的twist空间。

\section{为什么需要图形方法？}

\subsection{维度问题}

\begin{itemize}
\item \textbf{3D问题}：twist空间是6维，难以可视化
\item \textbf{平面问题}：twist空间是3维（$\omega_z, v_x, v_y$），可以可视化
\item 但即使是3维，直接绘制也不容易
\end{itemize}

\subsection{图形方法的优势}

\begin{itemize}
\item 将3维twist空间映射到2维平面上的点（CoR）
\item 可以在2D平面上直观地看到可行运动区域
\item 特别适用于静止接触约束的情况
\end{itemize}

\section{旋转中心（CoR）的概念}

\subsection{平面twist的表示}

对于平面问题，twist只有3个分量：
\begin{equation}
V = (\omega_z, v_x, v_y)
\label{eq:planar_twist}
\end{equation}

其中：
\begin{itemize}
\item $\omega_z$：绕$z$轴的角速度
\item $v_x, v_y$：在$xy$平面内的线速度
\end{itemize}

\subsection{旋转中心的定义}

\textbf{旋转中心}（Center of Rotation, CoR）是平面中在运动下保持\textbf{静止}的点。

\begin{definition}[旋转中心]
对于平面twist $V = (\omega_z, v_x, v_y)$，旋转中心的位置为：
\begin{equation}
\text{CoR} = \left(-\frac{v_y}{\omega_z}, \frac{v_x}{\omega_z}\right)
\label{eq:cor_position}
\end{equation}
\end{definition}

\textbf{特殊情况}：当$\omega_z = 0$时，CoR在\textbf{无穷远处}（纯平移运动）。

\subsection{旋转中心的物理意义}

\begin{itemize}
\item \textbf{旋转运动}：如果$\omega_z \neq 0$，物体绕CoR旋转
\item \textbf{平移运动}：如果$\omega_z = 0$，物体纯平移，CoR在无穷远
\item CoR是\textbf{螺旋轴}（screw axis）与平面的交点
\end{itemize}

\subsection{旋转方向的标记}

为了简化表示，我们只关心旋转的\textbf{方向}，而不关心速度大小：

\begin{itemize}
\item \textbf{'+'}：逆时针旋转（$\omega_z > 0$）
\item \textbf{'-'}：顺时针旋转（$\omega_z < 0$）
\item \textbf{'0'}：纯平移（$\omega_z = 0$）
\end{itemize}

\section{从Twist到CoR的映射}

\subsection{映射关系}

从平面twist $V = (\omega_z, v_x, v_y)$到CoR的映射：

\begin{equation}
\text{CoR}(V) = \begin{cases}
\left(-\frac{v_y}{\omega_z}, \frac{v_x}{\omega_z}\right) & \text{如果 } \omega_z \neq 0 \\
\text{无穷远（方向：}\frac{(v_x, v_y)}{|(v_x, v_y)|}\text{）} & \text{如果 } \omega_z = 0
\end{cases}
\label{eq:cor_mapping}
\end{equation}

\subsection{CoR空间的几何结构}

CoR空间由三部分组成：

\begin{enumerate}
\item \textbf{'+' CoR的平面}：所有逆时针旋转的CoR
\item \textbf{'-' CoR的平面}：所有顺时针旋转的CoR
\item \textbf{平移方向的圆}：所有纯平移的方向（在无穷远处）
\end{enumerate}

\subsection{映射的几何解释}

\begin{itemize}
\item 包含向量$V$的射线
\item 与$\omega_z = 1$平面相交 → '+' CoR
\item 与$\omega_z = -1$平面相交 → '-' CoR
\item 与$\omega_z = 0$圆相交 → 平移方向
\end{itemize}

\section{两个Twist的线性组合}

\subsection{线性组合的CoR}

给定两个不同的twist $V_1$和$V_2$，它们的线性组合：
\begin{equation}
V = k_1 V_1 + k_2 V_2, \quad k_1, k_2 \in \mathbb{R}
\label{eq:linear_combination}
\end{equation}

\textbf{关键性质}：线性组合的CoR位于通过$\text{CoR}(V_1)$和$\text{CoR}(V_2)$的\textbf{直线}上。

\subsection{两种情况}

\begin{enumerate}
\item \textbf{如果$\omega_{1z}$或$\omega_{2z}$非零}：
  \begin{itemize}
  \item 这条直线上的CoR可以有'+'或'-'标记
  \item 取决于$k_1$和$k_2$的符号
  \end{itemize}

\item \textbf{如果$\omega_{1z} = \omega_{2z} = 0$}：
  \begin{itemize}
  \item 两个twist都是纯平移
  \item 线性组合对应所有平移方向的集合
  \end{itemize}
\end{enumerate}

\section{正线性组合（Positive Linear Combination）}

\subsection{定义}

\textbf{正线性组合}（非负线性组合）：
\begin{equation}
V = \text{pos}(\{V_1, V_2\}) = \{k_1 V_1 + k_2 V_2 | k_1, k_2 \geq 0\}
\label{eq:positive_combination}
\end{equation}

这是一个以原点为根部的\textbf{平面twist锥}，$V_1$和$V_2$定义锥的边。

\subsection{情况1：相同符号的角速度}

如果$\omega_{1z}$和$\omega_{2z}$具有\textbf{相同的符号}：

\begin{itemize}
\item 正线性组合的CoR都具有该符号
\item CoR位于两个CoR之间的\textbf{线段}上
\item 例如：两个'+' CoR的正组合 → '+' CoR的线段
\end{itemize}

\begin{center}
\begin{tikzpicture}[scale=1.5]
% 两个'+' CoR
\filldraw[red] (0,0) circle (2pt) node[below left] {CoR$_1$ (+)};
\filldraw[red] (3,1) circle (2pt) node[above right] {CoR$_2$ (+)};

% 连接线段
\draw[thick,blue] (0,0) -- (3,1);
\fill[blue!20] (0,0) -- (3,1) -- (3.5,1.2) -- (0.5,0.2) -- cycle;

\node[blue] at (1.5,0.5) {正组合的CoR区域};
\node at (1.5,-0.5) {（线段）};
\end{tikzpicture}
\end{center}

\subsection{情况2：不同符号的角速度}

如果$\text{CoR}(V_1)$标记为'+'，$\text{CoR}(V_2)$标记为'-'：

\begin{itemize}
\item 正组合的CoR由包含两个CoR的\textbf{直线}组成
\item 但\textbf{减去}CoR之间的线段
\item 包括：
  \begin{itemize}
  \item 连接到CoR$_1$的'+' CoR射线
  \item 连接到CoR$_2$的'-' CoR射线
  \item 无穷远处的'0'点（对应平移）
  \end{itemize}
\end{itemize}

\begin{center}
\begin{tikzpicture}[scale=1.5]
% '+' CoR
\filldraw[red] (0,0) circle (2pt) node[below left] {CoR$_1$ (+)};

% '-' CoR
\filldraw[blue] (3,1) circle (2pt) node[above right] {CoR$_2$ (-)};

% 连接直线
\draw[thick,green] (-1,-0.33) -- (4,1.33);

% 标记中间不可行区域
\draw[thick,red,dashed] (0,0) -- (3,1);
\node[red] at (1.5,0.5) {不可行};

% 标记可行区域
\draw[thick,blue,->] (0,0) -- (-0.8,-0.27) node[left] {可行 (+)};
\draw[thick,green,->] (3,1) -- (3.8,1.27) node[right] {可行 (-)};

\node at (1.5,-0.5) {（直线，但中间段不可行）};
\end{tikzpicture}
\end{center}

\textbf{重要}：这个元素集合应该被视为\textbf{单个线段}（虽然通过无穷远），就像第一种情况一样。

\section{静止接触约束的CoR表示}

\subsection{基本思想}

对于\textbf{静止物体}和\textbf{可动物体}之间的接触，我们可以绘制不违反不可穿透约束的CoR。

\subsection{标记规则}

给定接触法线$\hat{n}$（指向可动物体内部）：

\begin{enumerate}
\item \textbf{接触法线上的点}：标记为'$\pm$'
\item \textbf{法线左侧的点}：标记为'+'（逆时针旋转可行）
\item \textbf{法线右侧的点}：标记为'-'（顺时针旋转可行）
\end{enumerate}

\subsection{物理解释}

\begin{itemize}
\item \textbf{'+'区域}：可以作为可动物体正角速度的CoR
\item \textbf{'-'区域}：可以作为可动物体负角速度的CoR
\item \textbf{'$\pm$'线}：接触法线，对应滚动或滑动接触
\end{itemize}

\subsection{接触标签的分配}

对于每个CoR，我们可以分配接触标签：

\begin{itemize}
\item \textbf{B（Breaking）}：分离接触
  \begin{itemize}
  \item CoR不在接触法线上
  \item 接触点正在远离
  \end{itemize}

\item \textbf{R（Rolling）}：滚动接触
  \begin{itemize}
  \item CoR在接触法线上（标记为'$\pm$'）
  \item 接触点处无相对运动
  \end{itemize}

\item \textbf{S（Sliding）}：滑动接触
  \begin{itemize}
  \item CoR在接触法线上（标记为'$\pm$'）
  \item 但接触点处有相对运动
  \item 进一步细分为：
    \begin{itemize}
    \item \textbf{$S_r$}：可动物体相对于固定约束向右滑动
    \item \textbf{$S_l$}：可动物体相对于固定约束向左滑动
    \end{itemize}
  \end{itemize}
\end{itemize}

\subsection{例子：单个静止接触}

\begin{center}
\begin{tikzpicture}[scale=2]
% 接触法线
\draw[thick,blue,->] (0,0) -- (0,2) node[above] {接触法线 $\hat{n}$};
\draw[thick,blue] (-0.1,0) -- (0.1,0) node[below] {接触点};

% 标记区域
\fill[red!20] (-2,0) -- (-2,2) -- (0,2) -- (0,0) -- cycle;
\node[red] at (-1,1) {+ 区域};

\fill[green!20] (0,0) -- (0,2) -- (2,2) -- (2,0) -- cycle;
\node[green] at (1,1) {- 区域};

% 法线上的标记
\draw[thick,blue] (0,0) -- (0,2);
\node[blue] at (0.3,1) {$\pm$};

% 标记一些CoR
\filldraw[red] (-1,0.5) circle (1.5pt) node[below] {+, B};
\filldraw[red] (-1,1.5) circle (1.5pt) node[above] {+, B};
\filldraw[blue] (0,1) circle (1.5pt) node[right] {$\pm$, R};
\filldraw[green] (1,0.5) circle (1.5pt) node[below] {-, B};
\filldraw[green] (1,1.5) circle (1.5pt) node[above] {-, B};
\end{tikzpicture}
\end{center}

\section{多个接触的情况}

\subsection{约束的并集}

如果有多个接触，我们只需取各个接触的约束和接触标签的\textbf{并集}。

\textbf{关键性质}：约束的这种并集意味着可行CoR区域是\textbf{凸的}，就像齐次多面体twist锥一样。

\subsection{例子：两个接触}

\begin{center}
\begin{tikzpicture}[scale=1.5]
% 第一个接触法线
\draw[thick,blue,->] (0,0) -- (0,2);
\fill[red!10] (-2,0) -- (-2,2) -- (0,2) -- (0,0) -- cycle;
\fill[green!10] (0,0) -- (0,2) -- (2,2) -- (2,0) -- cycle;

% 第二个接触法线
\draw[thick,red,->] (3,0) -- (4,1);
\fill[blue!10,rotate around={45:(3.5,0.5)}] (3.5,0.5) rectangle (5,2);

% 交集区域（可行区域）
\fill[purple!30] (0,0) -- (0,1.5) -- (1,2) -- (2,1.5) -- (2,0) -- cycle;
\node[purple] at (1,0.7) {可行CoR区域};

% 标记接触点
\filldraw[black] (0,0) circle (2pt) node[below left] {接触1};
\filldraw[black] (3,0) circle (2pt) node[below] {接触2};
\end{tikzpicture}
\end{center}

可行CoR区域是两个接触约束区域的\textbf{交集}。

\section{应用：形状闭合的可视化}

\subsection{形状闭合的CoR表示}

如果物体处于\textbf{形状闭合}：
\begin{itemize}
\item 可行twist集合只有原点（$V = 0$）
\item 对应的CoR：没有可行的CoR（除了特殊情况）
\item 所有方向都被约束阻止
\end{itemize}

\subsection{例子：四个接触固定物体}

\begin{itemize}
\item 四个静止接触
\item 每个接触贡献一个约束
\item 如果四个接触的约束区域没有交集（除了接触点），则物体被固定
\item 这是形状闭合的图形验证方法
\end{itemize}

\section{优势与局限性}

\subsection{优势}

\begin{enumerate}
\item \textbf{直观可视化}：在2D平面上看到3D twist空间
\item \textbf{几何直觉}：CoR的概念容易理解
\item \textbf{适用于静止约束}：特别适合分析静止接触
\end{enumerate}

\subsection{局限性}

\begin{enumerate}
\item \textbf{只适用于平面问题}：3D问题无法使用
\item \textbf{只适用于静止约束}：运动约束无法唯一表示
\item \textbf{特殊情况处理}：$\omega_z = 0$需要特殊处理
\end{enumerate}

\section{关键公式总结}

\subsection{CoR位置}

\begin{equation}
\text{CoR} = \left(-\frac{v_y}{\omega_z}, \frac{v_x}{\omega_z}\right), \quad \omega_z \neq 0
\end{equation}

\subsection{正线性组合}

\begin{equation}
\text{pos}(\{V_1, V_2\}) = \{k_1 V_1 + k_2 V_2 | k_1, k_2 \geq 0\}
\end{equation}

\subsection{标记规则}

\begin{itemize}
\item 法线左侧：'+'
\item 法线右侧：'-'
\item 法线上：'$\pm$'
\end{itemize}

\section{总结}

12.1.6章节的核心内容：

\begin{enumerate}
\item \textbf{CoR表示}：将平面twist映射到2D平面上的点

\item \textbf{线性组合}：两个twist的线性组合对应CoR的直线

\item \textbf{正线性组合}：形成CoR的线段或射线

\item \textbf{静止约束}：可以用CoR区域表示可行运动

\item \textbf{接触标签}：每个CoR可以分配接触标签（B、R、S）

\item \textbf{多个接触}：约束区域的交集给出可行CoR区域
\end{enumerate}

这种方法为理解平面接触运动学提供了强大的可视化工具。

\end{document}

